% =====================================================================
%  CHAPTER 21: 社会病防治
% =====================================================================
\chapter{社会病防治}
\setcounter{chapter}{21}
\chapterintro{
掌握社会病的概念、特点及人际暴力的相关理论（\textbf{教学难点}）；熟悉自杀的概念、分类、分布特征和社会根源；了解攻击、暴力、非故意伤害、成瘾行为、吸毒以及与性行为相关的社会病等问题。
}

% ----- 21.1 社会病概述 -----
\section{社会病概述}

\begin{defbox}
\textbf{社会病（Sociopathy / Social Pathology）}：由于社会原因（社会制度、社会结构、社会文化、社会转型等）而产生的\keyword{损害社会功能和人群健康}的社会问题。社会病既是社会问题（社会层面的功能失调），也是健康问题（对人群身心健康造成严重损害）。

\textbf{辨析：社会病与社会问题、越轨行为}
\begin{itemize}[leftmargin=*]
    \item \imp{社会问题（Social Problem）}：广义上任何影响社会正常运行的问题，社会病是社会问题中对人群健康影响最直接的一类
    \item \imp{越轨行为（Deviant Behavior）}：偏离或违反社会规范的行为，多数社会病涉及越轨行为，但并非所有越轨行为都是社会病
    \item \imp{反社会人格障碍（Antisocial Personality Disorder）}：精神医学诊断，主要表现为长期漠视和侵犯他人权利的行为模式，与社会病概念不同但有一定交叉
\end{itemize}
\end{defbox}

\subsection{社会病的特点}

\begin{keybox}
\textbf{社会病的五大特点：}
\begin{enumerate}[leftmargin=*]
    \item \imp{社会根源性}——社会病的根源在于社会因素（社会不平等、制度缺陷、文化冲突等），而非纯粹的个体生物学问题或个体道德问题
    \item \imp{群体性/流行性}——社会病涉及一定规模的人群，具有一定的流行特征
    \item \imp{严重危害性}——对社会公共安全、社会秩序、经济发展和人群身心健康造成严重危害
    \item \imp{复杂性和综合性}——社会病的成因复杂多元，涉及多个层面：个体因素与宏观社会因素相互交织，往往与贫困、教育程度低、社会歧视等社会不利条件相关
    \item \imp{可预防性}——可以通过社会政策、法制建设、教育宣传等社会干预手段有效预防和控制
\end{enumerate}
\end{keybox}

% ----- 21.2 攻击与暴力 -----
\section{攻击与暴力}

\subsection{基本概念}

\begin{defbox}
\textbf{攻击（Aggression）}：以伤害他人或损害物体为目的的行为。可以分为：工具性攻击（以达成某种目标为目的）和敌意性攻击（以伤害他人本身为目的）。

\textbf{暴力（Violence）}：故意运用躯体力量或权力，对自身、他人、群体或社会进行威胁或伤害，造成或极有可能造成损伤、死亡、心理伤害、发育不良或剥夺的行为。
\end{defbox}

\subsection{人际暴力（教学难点*）}

\begin{keybox}
\imp{（教学难点*）}\textbf{人际暴力（Interpersonal Violence）：}

\textbf{WHO将暴力分为三大类：}
\begin{enumerate}[leftmargin=*]
    \item \imp{自我指向暴力（Self-directed Violence）}：自杀行为（自杀、自杀未遂、自杀意念）、自伤行为
    \item \imp{人际暴力（Interpersonal Violence）}：家庭暴力/亲密伴侣暴力（配偶虐待、儿童虐待、老年人虐待）、社区暴力（熟人暴力和陌生人暴力）
    \item \imp{集体暴力（Collective Violence）}：社会暴力、政治暴力、经济暴力
\end{enumerate}

\textbf{家庭暴力（Domestic Violence）的特点：}
\begin{itemize}[leftmargin=*]
    \item 隐蔽性强——发生在家庭内部，不易为外人所知
    \item 反复性和循环性——暴力往往呈周期性循环模式
    \item 多重伤害——造成躯体、心理和社会功能的全面损害
    \item 代际传递——暴力家庭中成长的儿童成年后更可能成为施暴者或受害者
\end{itemize}

\textbf{暴力的社会根源：}
\begin{enumerate}[leftmargin=*]
    \item 社会不平等——贫富差距、社会排斥
    \item 社会文化因素——暴力亚文化、"荣誉"文化、性别不平等
    \item 社会解组——传统社会规范和约束机制失效
    \item 酒精和药物滥用——降低行为控制力
    \item 儿童期不良经历（ACEs）——遭受虐待或忽视
\end{enumerate}
\end{keybox}

% ----- 21.3 自杀 -----
\section{自杀}

\begin{defbox}
\textbf{自杀（Suicide）}：个体在意识清晰的情况下，自愿地采取某种手段结束自己生命的行为。

\textbf{自杀行为的分类：}
\begin{enumerate}[leftmargin=*]
    \item \imp{自杀意念（Suicidal Ideation）}——有结束生命的想法但未付诸行动
    \item \imp{自杀未遂（Attempted Suicide）}——实施了自杀行为但未导致死亡
    \item \imp{自杀死亡（Completed Suicide）}——自杀行为导致死亡
\end{enumerate}
\end{defbox}

\subsection{自杀的发生率与分布特征}

\begin{keybox}
\begin{itemize}[leftmargin=*]
    \item 全球：每年约80万人死于自杀（每40秒1人），自杀未遂人数约为自杀死亡人数的20倍
    \item 中国：曾经是全球自杀率最高的国家之一（1990年代年自杀率约23/10万），近年来显著下降（约7/10万），这与社会经济发展和农药管理改善有关
    \item 中国自杀的独特模式：①\imp{农村显著高于城市}（3-5倍，主要与农药可及性高有关）；②老年人为最高风险人群；③中国女性自杀率曾高于男性（全球大多数国家男性高于女性），近年已趋近国际模式
    \item 全球：自杀死亡以男性为主（男性:女性 $\approx$ 3:1），但自杀未遂以女性居多
    \item 高危人群：老年人、青少年（15-29岁第二大死因）、精神障碍患者、LGBTQ群体
\end{itemize}
\end{keybox}

\subsection{自杀的社会根源}

\begin{enumerate}[leftmargin=*]
    \item \imp{社会整合不足}——Durkheim的四种自杀类型：
    \begin{itemize}
        \item \imp{利己型自杀（Egoistic）}：个人与社会整合不足，个人感到孤立无援、生存无意义
        \item \imp{利他型自杀（Altruistic）}：个人过度整合于社会，为群体利益主动牺牲生命（如军人为国捐躯、宗教殉道）
        \item \imp{失范型自杀（Anomic）}：社会规范突然崩溃或缺失（如经济危机、突然暴富或破产），使个人行为失去指引
        \item \imp{宿命型自杀（Fatalistic）}：社会规范过度压抑个人，个人感到绝望、无处可逃（如奴隶、囚犯）
    \end{itemize}
    \item \imp{精神障碍}——抑郁症、物质成瘾、精神分裂症等是最重要的危险因素
    \item \imp{社会孤立}——缺乏社会联系和支持网络
    \item \imp{经济压力}——失业、债务、贫困等经济因素
    \item \imp{社会文化因素}——某些文化中自杀被视为"光荣"行为
    \item \imp{媒体效应}——不当的自杀报道可引发模仿自杀（Werther效应）
\end{enumerate}

\subsection{自杀的预防}

\begin{enumerate}[leftmargin=*]
    \item 限制自杀手段的可及性——农药管理、枪支管控、高层建筑和高桥防护
    \item 早期识别和干预——高危人群筛查和跟踪；危机干预热线（24小时心理援助热线）
    \item 加强精神卫生服务——提高抑郁症等精神障碍的识别率和治疗率
    \item 负责任的媒体报道——遵循WHO自杀报道指南，减少Werther效应
    \item 加强社会支持——社区关怀、家庭支持、社会再融入
    \item 综合性国家自杀预防战略——WHO建议各国制定国家层面的自杀预防战略
\end{enumerate}

% ----- 21.4 非故意伤害 -----
\section{非故意伤害}

\begin{defbox}
\textbf{非故意伤害（Unintentional Injury）}：由于非故意的意外事件造成的躯体伤害，包括交通事故、跌倒、溺水、中毒、烧伤等。非故意伤害是全球重要的公共卫生问题，是儿童和青少年死亡的重要原因。
\end{defbox}

\textbf{主要类型：}交通事故伤害（全球第八大死因，每年约135万人死亡）、跌倒（老年人最主要伤害原因）、溺水（1-14岁儿童首位伤害死因）、中毒、烧伤和烫伤。

% ----- 21.5 成瘾行为 -----
\section{成瘾行为}

\subsection{吸毒}

\begin{keybox}
\textbf{吸毒（Drug Abuse/Addiction）的健康和社会危害：}
\begin{enumerate}[leftmargin=*]
    \item 对个体的危害——躯体健康损害（器官衰竭、药物过量死亡）、精神障碍（精神病性症状）、HIV/AIDS和肝炎通过共用针具传播、社会功能丧失
    \item 对家庭和社会的危害——家庭破裂和经济崩溃、药物相关的犯罪（暴力、盗窃）、社会生产力的巨大损失、社会治理成本急剧增加
\end{enumerate}

\textbf{吸毒的社会防治：}
\begin{enumerate}[leftmargin=*]
    \item 法律和执法——严厉打击毒品贩运和生产
    \item 预防教育——学校、社区和家庭的毒品预防教育
    \item 减害策略（Harm Reduction）——美沙酮维持治疗（MMT）、针具交换项目（NSP）
    \item 治疗和康复——社区戒毒、强制隔离戒毒、自愿戒毒，重点是回归社会
\end{enumerate}
\end{keybox}

\subsection{问题饮酒行为}

\begin{defbox}
\textbf{问题饮酒（Problem Drinking）}：饮酒导致健康损害、社会功能受损或法律问题的饮酒模式。包括酒精滥用（Alcohol Abuse）和酒精依赖（Alcohol Dependence/Alcoholism）。
\end{defbox}

\textbf{过量饮酒的健康危害：}肝脏疾病（酒精性肝炎、肝硬化、肝癌）、心血管疾病、多种癌症（口腔、咽喉、食管、肝癌等）、神经系统损害、胎儿酒精综合征。

% ----- 21.6 与性行为相关的社会病 -----
\section{与性行为相关的社会病}

\begin{enumerate}[leftmargin=*]
    \item \imp{性传播疾病（STDs）}——艾滋病（AIDS）、梅毒、淋病、尖锐湿疣等。全球约3800万HIV/AIDS感染者；性传播是我国HIV的主要传播途径
    \item \imp{青少年妊娠（Teenage Pregnancy）}——青少年因身心发育未成熟和缺乏社会支持，面临更高的孕产期并发症风险，同时造成教育中断和社会经济劣势的代际循环
    \item \imp{性暴力}——严重侵犯人权的社会问题，给受害者造成长期的身心创伤
\end{enumerate}

\textbf{社会防治：}综合性性教育、提高避孕意识和可及性、安全套推广（针对于STDs和HIV）、反对性暴力和加强社会保护。

% ----- 21.7 精神障碍 -----
\section{精神障碍}

\begin{defbox}
\textbf{精神障碍（Mental Disorders）}：以认知、情绪、行为异常为主要表现的一大类精神疾病的总称。

精神障碍不仅是医学问题，也是重要的\keyword{公共卫生和社会问题}——全球约10亿人患有精神障碍，精神障碍占全球DALY的7\%以上。由于病耻感（STIGMA），大多数患者未能获得适当治疗。
\end{defbox}

% ----- 复习题 -----
\section{复习题}

\subsection*{一、选择题}
\begin{enumerate}[leftmargin=*, itemsep=3pt]
    \item 社会病的首要特点是：
    \begin{enumerate}
        \item[(A)] 生物因素为主
        \item[(B)] 社会根源性
        \item[(C)] 不可预防性
        \item[(D)] 个体道德问题
    \end{enumerate}
    \item WHO将暴力分为三大类，不包括：
    \begin{enumerate}
        \item[(A)] 自我指向暴力
        \item[(B)] 人际暴力
        \item[(C)] 财产暴力
        \item[(D)] 集体暴力
    \end{enumerate}
    \item Durkheim将自杀分为四种类型，其中社会整合不足导致的是：
    \begin{enumerate}
        \item[(A)] 利己型自杀
        \item[(B)] 利他型自杀
        \item[(C)] 失范型自杀
        \item[(D)] 宿命型自杀
    \end{enumerate}
    \item 减害策略（Harm Reduction）不包括以下哪项？
    \begin{enumerate}
        \item[(A)] 美沙酮维持治疗
        \item[(B)] 针具交换项目
        \item[(C)] 鼓励吸毒
        \item[(D)] 安全套推广
    \end{enumerate}
    \item 以下哪个群体是自杀风险最高的年龄段之一？
    \begin{enumerate}
        \item[(A)] 5-10岁儿童
        \item[(B)] 15-29岁青少年/青年（第二大死因）
        \item[(C)] 30-40岁中年人
        \item[(D)] 50-60岁中年人
    \end{enumerate}
\end{enumerate}

\subsection*{二、名词解释}
\begin{enumerate}[leftmargin=*, itemsep=3pt]
    \item \textbf{社会病（Social Pathology）}
    \item \textbf{自杀（Suicide）}
    \item \textbf{越轨行为（Deviant Behavior）}
    \item \textbf{减害策略（Harm Reduction）}
\end{enumerate}

\subsection*{三、简答题}
\begin{enumerate}[leftmargin=*, itemsep=3pt]
    \item 简述社会病的特点。
    \item 简述WHO对暴力的三大分类及人际暴力的主要类型。
    \item 简述自杀的社会根源及预防策略。
    \item 简述吸毒对个体和社会的主要危害。
\end{enumerate}

\subsection*{参考答案要点}

\subsubsection*{一、选择题}
1. (B)；2. (C)；3. (A)；4. (C)；5. (B)

\subsubsection*{二、名词解释}
\begin{enumerate}[leftmargin=*]
    \item \textbf{社会病}：由于社会原因产生的损害社会功能和人群健康的社会问题，具有社会根源性、群体性、严重危害性、复杂性和可预防性等特点。
    \item \textbf{自杀}：个体在意识清晰的情况下自愿采取某种手段结束自己生命的行为，包括自杀意念、自杀未遂和自杀死亡。
    \item \textbf{越轨行为}：偏离或违反社会规范的行为，多数社会病涉及越轨行为，但并非所有越轨行为都是社会病。
    \item \textbf{减害策略}：以减少药物滥用负面后果（而非强制禁绝）为目标的公共卫生策略，如美沙酮维持治疗、针具交换项目。
\end{enumerate}

\subsubsection*{三、简答题}
\begin{enumerate}[leftmargin=*]
    \item 五大特点：社会根源性（根源于社会不平等和制度缺陷）、群体性（涉及较大规模人群）、严重危害性（公共安全和身心健康）、复杂性和综合性（多层面因素交织、与社会不利条件相关）、可预防性（社会政策、法制、教育等措施有效）。
    \item 三大类：①自我指向暴力（自杀、自伤）；②人际暴力（家庭暴力/亲密伴侣暴力、社区暴力）；③集体暴力（社会、政治、经济暴力）。人际暴力主要类型：家庭暴力（配偶、儿童、老年人虐待）和社区暴力（熟人暴力和陌生人暴力）。
    \item 社会根源：社会整合不足（Durkheim理论）、精神障碍、社会孤立、经济压力、文化因素、媒体效应（Werther效应）。预防策略：限制自杀手段可及性、高危人群筛查和危机干预、加强精神卫生服务、负责任媒体报道、加强社会支持、制定国家自杀预防战略。
    \item 个体危害：躯体健康损害（器官衰竭、药物过量死亡）、精神障碍、传染病传播（共用针具→HIV/AIDS/肝炎）、社会功能丧失。社会危害：家庭破裂与经济崩溃、药物相关犯罪、社会生产力损失、治理成本增加。
\end{enumerate}

\newpage
